\section{Related work}\label{sec:related}

\dbsp using non-nested streams is a simplified instance of a Kahn 
network~\cite{kahn-ifip74}.  Johnson~\cite{johnson-phd83}
studies a very similar computational model without nested streams and its 
expressiveness. The implementation of such streaming models of computation and their
relationship to dataflow machines has been studied by Lee~\cite{lee-ieee95}.
Lee~\cite{lee-ifip93} also introduced streams of streams and the $\lift{\zm}$ operator.

In \secref{sec:extensions} we have discussed the connections with window and stream database 
queries~\cite{arasu-tr02,aurora}.

Incremental view maintenance (e.g.~\cite{gupta-idb93}) is
surveyed in~\cite{gupta-idb95}; a large bibliography is present in~\cite{motik-ai19}. 
Its most formal aspect is propagating ``deltas'' through algebraic expressions:
$Q(R+\Delta R)=Q(R)+\Delta Q(R,\Delta R)$. This work eventually crystallized in~\cite{koch-pods16}. DBSP 
incrementalization is both more modular and more fine-grain since it deals with streams of updates. 
Both~\cite{koch-pods10} and~\cite{green-tcs11} use \zrs to uniformly model insertions/deletions.

Picallo et al.~\cite{picallo-scop19} provide a general solution to IVM for
rich languages.  \dbsp requires a group structure on the values operated on; 
this assumption has two major practical benefits: it simplifies the mathematics considerably
(e.g., Picallo uses monoid actions to model changes), and it provides a general, simple
algorithm (\ref{algorithm-inc}) for incrementalizing arbitrary programs.  The downside of 
\dbsp is that one has to find a suitable group structure (e.g., \zrs for sets) to ``embed'' 
the computation.  Picallo's notion of ``derivative'' is not unique: they need creativity to choose
the right derivative definition, we need creativity to find the right group structure.

\begin{comment}
The main problem that change structures address is that the types used in programs are not
closed under subtraction (e.g., the delta between two sets is not a set). 
Although a relational \dbsp circuit computes
only on positive \zr values, its incremental version may compute on negative 
values, but the equivalence of the two programs guarantees correctness even though the
type system of \zrs does not.  \val{Safe to delete this para}
\end{comment}

Many heuristic algorithms were published for Datalog-like languages, e.g., 
counting based approaches~\cite{Dewan-iis92,motik-aaai15} that maintain the number of derivations,
DRed~\cite{gupta-sigmod93} and its variants~\cite{Ceri-VLDB91,Wolfson-sigmod91,%
Staudt-vldb96,Kotowski-rr11,Lu-sigmod95,Apt-sigmod87}, the backward-forward algorithm and variants~\cite{motik-aaai15,
Harrison-wdd92,motik-ai19}.  \dbsp is more general than these approaches.  
Interestingly, the \zrs multiplicities in our relational implementation 
are related to the counting-number-of-derivations approaches.

\dbsp is tightly related to Differential Dataflow (DD)~\cite{mcsherry-cidr13, murray-sosp13}
and its theoretical foundations~\cite{abadi-fossacs15} (and recently~\cite{mchserry-vldb20,chothia-vldb16}).
All \dbsp operators are based on DD operators.  DD's computational model is more powerful than
\dbsp, since it allows past values in a stream to be "updated".
In contrast, our model assumes that the inputs of a computation arrive in the time order while allowing 
for nested time domains via the modular lifting transformer.  However, \dbsp can express both
incremental and non-incremental computations; in essence \dbsp is ``deconstructing'' DD into 
simple component building blocks; the core Proposition~\ref{prop-inc-properties} and
the Algorithm based on it~\ref{algorithm-inc} are new contributions.
